% ============================================================
% INFORMAÇÕES DO PROJETO - UNESP
% ============================================================

% ----- DADOS DO ALUNO (EDITAR) -----
\newcommand{\autor}{Nome do Aluno}
\newcommand{\autorgenero}{f} % m | f (Define se é "Graduando" ou "Graduanda")

% ----- TIPO DE CURSO (EDITAR) -----
% Descomente a sua escolha:
%\newcommand{\tipocurso}{licenciatura}
\newcommand{\tipocurso}{bacharelado} 

% ----- DADOS DO ORIENTADOR (EDITAR) -----
\newcommand{\orientador}{Nome do Orientador}
\newcommand{\orientadortitulo}{Profa. Dra.} % Ou Prof. Dr.
\newcommand{\orientadorgenero}{f} % m | f

% ----- DADOS DO COORIENTADOR (OPCIONAL) -----
% Deixe o nome vazio se não houver coorientador: \newcommand{\coorientador}{}
\newcommand{\coorientador}{} 
\newcommand{\coorientadortitulo}{Prof. Dr.} % Ou Profa. Dra.
\newcommand{\coorientadorgenero}{m} % m | f

% ----- CAMPUS E DATA (EDITAR) -----
\newcommand{\cidade}{Rio Claro}
\newcommand{\estado}{SP}
\newcommand{\campus}{\textit{Câmpus} de \cidade}
\newcommand{\ano}{\the\year} 
\newcommand{\mes}{\ifcase\month\or Janeiro\or Fevereiro\or Março\or Abril\or Maio\or Junho\or Julho\or Agosto\or Setembro\or Outubro\or Novembro\or Dezembro\fi}

% ----- TÍTULO E SUBTÍTULO (EDITAR) -----
\newcommand{\titulo}{Título do Projeto de TCC:}
\newcommand{\subtitulo}{Subtítulo do Projeto (se houver)}

% ============================================================
% CONTROLE DOS PRÉ-TEXTUAIS E PÓS-TEXTUAIS (EDITAR)
% ============================================================
\newcommand{\ativarApêndices}{0}
\newcommand{\ativarAnexos}{0}

% ============================================================
% CONFIGURAÇÕES INSTITUCIONAIS (NÃO ALTERAR)
% ============================================================

\newcommand{\universidade}{Universidade Estadual Paulista}
\newcommand{\universidadenome}{Júlio de Mesquita Filho}
\newcommand{\universidadeabreviada}{UNESP}

% --- SELEÇÃO DO CURSO (ESCOLHA APENAS UM, DESCOMENTANDO A LINHA) ---

% === INSTITUTO DE BIOCIÊNCIAS (IB) ===
 \def\meuCurso{ciencias-biologicas}
% \def\meuCurso{ecologia}
% \def\meuCurso{educacao-fisica}
% \def\meuCurso{pedagogia}

% === INSTITUTO DE GEOCIÊNCIAS (IGCE) ===
% \def\meuCurso{geografia}
% \def\meuCurso{ciencias-computacao}
% \def\meuCurso{engenharia-ambiental}
% \def\meuCurso{fisica}
% \def\meuCurso{geologia}
% \def\meuCurso{matematica}

% --- MAPEAMENTO CURSO -> INSTITUTO, DEPARTAMENTO, SIGLA ---
\IfStrEq{\meuCurso}{ciencias-biologicas}{%
    \newcommand{\institutenome}{Instituto de Biociências}
    \newcommand{\institutoabreviado}{IB}
    \newcommand{\curso}{Ciências Biológicas}
    \newcommand{\departamento}{Departamento de Biologia Geral e Aplicada}
}{%
\IfStrEq{\meuCurso}{ecologia}{%
    \newcommand{\institutenome}{Instituto de Biociências}
    \newcommand{\institutoabreviado}{IB}
    \newcommand{\curso}{Ecologia}
    \newcommand{\departamento}{Departamento de Biodiversidade}
}{%
\IfStrEq{\meuCurso}{educacao-fisica}{%
    \newcommand{\institutenome}{Instituto de Biociências}
    \newcommand{\institutoabreviado}{IB}
    \newcommand{\curso}{Educação Física}
    \newcommand{\departamento}{Departamento de Educação}
}{%
\IfStrEq{\meuCurso}{pedagogia}{%
    \newcommand{\institutenome}{Instituto de Biociências}
    \newcommand{\institutoabreviado}{IB}
    \newcommand{\curso}{Pedagogia}
    \newcommand{\departamento}{Departamento de Educação}
}{%
\IfStrEq{\meuCurso}{ciencias-computacao}{%
    \newcommand{\institutenome}{Instituto de Geociências e Ciências Exatas}
    \newcommand{\institutoabreviado}{IGCE}
    \newcommand{\curso}{Ciências da Computação}
    \newcommand{\departamento}{Departamento de Estatística, Matemática Aplicada e Computação}
}{%
\IfStrEq{\meuCurso}{engenharia-ambiental}{%
    \newcommand{\institutenome}{Instituto de Geociências e Ciências Exatas}
    \newcommand{\institutoabreviado}{IGCE}
    \newcommand{\curso}{Engenharia Ambiental}
    \newcommand{\departamento}{Departamento de Geologia Aplicada e Ambiental}
}{%
\IfStrEq{\meuCurso}{fisica}{%
    \newcommand{\institutenome}{Instituto de Geociências e Ciências Exatas}
    \newcommand{\institutoabreviado}{IGCE}
    \newcommand{\curso}{Física}
    \newcommand{\departamento}{Departamento de Física}
}{%
\IfStrEq{\meuCurso}{geografia}{%
    \newcommand{\institutenome}{Instituto de Geociências e Ciências Exatas}
    \newcommand{\institutoabreviado}{IGCE}
    \newcommand{\curso}{Geografia}
    \newcommand{\departamento}{Departamento de Geografia e Planejamento Ambiental}
}{%
\IfStrEq{\meuCurso}{geologia}{%
    \newcommand{\institutenome}{Instituto de Geociências e Ciências Exatas}
    \newcommand{\institutoabreviado}{IGCE}
    \newcommand{\curso}{Geologia}
    \newcommand{\departamento}{Departamento de Geologia}
}{%
\IfStrEq{\meuCurso}{matematica}{%
    \newcommand{\institutenome}{Instituto de Geociências e Ciências Exatas}
    \newcommand{\institutoabreviado}{IGCE}
    \newcommand{\curso}{Matemática}
    \newcommand{\departamento}{Departamento de Matemática}
}{%
    % --- FALLBACK ---
    \newcommand{\institutenome}{Instituto}
    \newcommand{\institutoabreviado}{-}
    \newcommand{\curso}{Curso}
    \newcommand{\departamento}{Departamento}
}}}}}}}}}}

% ============================================================
% COMANDOS DERIVADOS (NÃO ALTERAR - LÓGICA AUTOMÁTICA)
% ============================================================

% ----- FORMATAÇÃO DO GÊNERO DO ALUNO -----
\ifthenelse{\equal{\autorgenero}{m}}{%
    \newcommand{\grauautor}{Graduando}
}{%
    \newcommand{\grauautor}{Graduanda}
}

% ----- FORMATAÇÃO DO GÊNERO DO ORIENTADOR -----
\ifthenelse{\equal{\orientadorgenero}{m}}{%
    \newcommand{\orientadorfuncao}{Orientador}
    \newcommand{\orientadorpreposicao}{pelo} % Para a frase da folha de rosto
}{%
    \newcommand{\orientadorfuncao}{Orientadora}
    \newcommand{\orientadorpreposicao}{pela} % Para a frase da folha de rosto
}
\newcommand{\orientadornome}{\orientadortitulo\ \orientador}

% ----- FORMATAÇÃO DO GÊNERO DO COORIENTADOR -----
\ifthenelse{\equal{\coorientador}{}}{%
    \newcommand{\temcoorientador}{false}
    \newcommand{\coorientadornome}{}
}{%
    \newcommand{\temcoorientador}{true}
    \ifthenelse{\equal{\coorientadorgenero}{m}}{%
        \newcommand{\coorientadorfuncao}{Coorientador}
    }{%
        \newcommand{\coorientadorfuncao}{Coorientadora}
    }
    \newcommand{\coorientadornome}{\coorientadortitulo\ \coorientador}
}

% ----- TEXTO DA NATUREZA DO TRABALHO (FOLHA DE ROSTO) -----
\newcommand{\naturezaTrabalho}{%
    Projeto de Trabalho de Conclusão de Curso de \tipocurso\ em %
    \curso, orientado \orientadorpreposicao\ \orientadornome, %
    apresentado ao Conselho do Curso de Graduação em \curso\ %
    da \universidade\ \enquote{\universidadenome} -- %
    UNESP / \campus.%
}

% ============================================================
% FIM DO ARQUIVO infos.tex
% ============================================================